% 第8章 8.3.1、8.3.3：Lipschitz 类的网格量化覆盖。L=1，eps=0.4，delta=eps/(4L)=0.1，eta=eps/2=0.2。
% 空间网格 x_i=0.1,0.3,0.5,0.7,0.9 覆盖 [0,1]（每点到最近网格点不超过 delta，竖虚线）；值域按 eta 分层（横灰线）。
% f(x)=0.1+0.3 sin(3x)（蓝），g 为过 (0,0),(0.1,0.05),(0.3,0.22),(0.5,0.25),(0.7,0.39),(0.9,0.38),(1,0.3) 的折线（绿，斜率绝对值<=0.85）；
% 二者在每个网格点落在同一层（同一记录），一致距离 <= 2L delta+eta=0.4。
\begin{tikzpicture}[font=\figurefont\footnotesize]
\pgfplotsset{every axis/.style={axis lines=left,
  x axis line style={draw=FigureStroke,line width=.6pt},y axis line style={draw=FigureStroke,line width=.6pt},
  tick style={draw=FigureStroke},tick label style={font=\figurefont\scriptsize,text=FigureInk},
  title style={font=\figurefont\footnotesize},label style={font=\figurefont\scriptsize},
  legend style={draw=none,fill=none,font=\figurefont\scriptsize,cells={anchor=west}}}}
\begin{axis}[width=74mm,height=44mm,xmin=0,xmax=1.03,ymin=0,ymax=0.6,xtick={0.1,0.3,0.5,0.7,0.9},
  xticklabels={$x_1$,$x_2$,$x_3$,$x_4$,$x_5$},ytick={0,0.2,0.4,0.6},
  legend style={at={(1.03,0.5)},anchor=west}]
\foreach \y in {0.2,0.4}{\edef\tmp{\noexpand\draw[FigureGrid,line width=.6pt] (axis cs:0,\y)--(axis cs:1.03,\y);}\tmp}
\foreach \x in {0.1,0.3,0.5,0.7,0.9}{\edef\tmp{\noexpand\draw[FigureMuted,line width=.4pt,dashed] (axis cs:\x,0)--(axis cs:\x,0.6);}\tmp}
\addplot[draw=FigureBlue,line width=1.1pt,domain=0:1,samples=120] {0.1+0.3*sin(deg(3*x))};
\addplot[draw=FigureGreen,line width=1.1pt,dashed] coordinates {(0,0) (0.1,0.05) (0.3,0.22) (0.5,0.25) (0.7,0.39) (0.9,0.38) (1,0.3)};
\addplot[only marks,mark=*,mark size=1.4pt,FigureBlue] coordinates {(0.1,0.18867) (0.3,0.33500) (0.5,0.39925) (0.7,0.35926) (0.9,0.22821)};
\addplot[only marks,mark=square*,mark size=1.3pt,FigureGreen] coordinates {(0.1,0.05) (0.3,0.22) (0.5,0.25) (0.7,0.39) (0.9,0.38)};
\legend{$f$,$g$}
\draw[<->,FigureInk,line width=.5pt] (axis cs:1.0,0.4)--node[left,font=\figurefont\scriptsize]{$\eta$}(axis cs:1.0,0.6);
\end{axis}
\end{tikzpicture}
